\section{Conclusion}

We characterised context-dependent selection among four established urban
routing policies in \NumRuns{} microsimulation runs over \NumContexts{}
controlled contexts, with every policy decision made from information available
before activation and every comparison paired under common random numbers.

The portfolio is genuinely complementary: three of the four policies are
strictly best in some context by a margin the seeds resolve, the winner map
depends on more than one factor, and \NumParetoNonSingleton{} of contexts have a
non-singleton Pareto set across journey time, \CO{} and stopped delay. The
reliability-aware policy, which almost never minimises journey time, minimises
stopped delay in \NumSpecCThreePct{} of contexts. Complementarity in this
portfolio lives in the criterion vector, not in the scalar cost.

The value of exploiting it is nevertheless \NumHeadroom~s of system mean journey
time per vehicle --- \NumHeadroomPct{} --- because the best fixed policy is
near-dominant, with an upside \NumAsymRatio$\times$ its downside. Against this,
choosing the wrong fixed policy costs \NumFixedWorstPenaltyPct{}.

Within that small margin the methodological findings are clear. A selector
trained to predict the winner's identity is more accurate than the fixed policy
and more expensive than it; a selector trained to predict each policy's
advantage over the fixed default is less accurate and cheaper, closing
\NumBFourGapClosed{} of the available gap. Adaptation helps in-distribution and
harms under every extrapolation and shift we constructed, and a support-and-
confidence gate prevented all of that harm by declining to act --- correctly,
because the uncertainty exceeds the effect.

The operational conclusion is not that adaptive routing-policy selection cannot
work. It is that in this environment the decision worth making is which fixed
policy to run and which criterion to optimise, and that the headroom for
per-context adaptation can be measured, cheaply and before any model is built,
by comparing the best fixed policy against the retrospective best and both
against the noise floor of the same comparison.
